\begin{afterword}
\summaryhead

The post asks why a discovery should be worth less because someone else already knows the answer. Being the first to solve a major mystery is said to be a tremendous high; the physicist who first knew why the stars shine is the example. The author weighs three readings of that joy. The most charitable is that a long struggle gives a deep understanding that falls into place at once. A less charitable one is effort justification, valuing a thing because it cost more. The least charitable is status.

If the joy requires being first, the author argues, it is scarce, and in a universe large enough to contain aliens, alternate Newtons or distant duplicates, someone has always been first. The author calls this a reductio. The consistent solution is to treat mystery as personal: if you do not know the answer, it is a mystery to you, whoever else knows. Telling one's civilization something new remains a separate reward. The post ends by inviting the reader to see how much of the world is still, to them, unknown.

\responsehead

The post's conclusion is modest and well argued: the pleasure of finding something out does not depend on no one else knowing it. Its three readings of the joy of first discovery are presented as readings, from charitable to uncharitable, and the post grants that the reward of telling one's civilization something new is real.

The evidence it describes is less exact than its argument. For effort justification it gives two studies without citing them. The classic study of initiations, by Aronson and Mills, used women joining a discussion group, not fraternity pledges. The wine result does not illustrate effort at all: tasters told a higher price reported more pleasure, and a price label is not a cost the taster bears. The claim that the thrill of first discovery is ``attested by numerous sources'' names none. None of this harms the conclusion, which does not depend on these studies.

The post misuses a physics term. The post's ``Standard Model of physics'' is the theory of particles, and says nothing about infinite space, inflation or branching worlds. The three claims are Tegmark's first three levels of parallel universes; Tegmark treats infinite space as part of the standard cosmological model, and the other two as following from chaotic inflation and from unitary quantum mechanics, the many-worlds reading. The argument needs only infinite space.

The reductio, the argument that the view leads to an absurd result, covers only one version of the view. It shows that joy which requires no one anywhere to know is never available. The opening story, and the Nobel Prize paragraph near the end, concern being first on one Earth, and other Earths make no difference to that. For that version the post's case is the earlier one: exclusivity makes the joy scarce.

In short: a modest and sound conclusion, that the joy of finding out is personal, reached partly through a reductio that applies only to the view that no one anywhere may know.
\end{afterword}
